\section{Branching}

Branching is one of the most fundamental patterns of behaviour observable in
living systems. It is a basic means for the ordered structures of living
entities to expand indefinitely in response to resource supply. Dividing one
individual into two allows biomass, and hence dissipative potential, to
increase without paying the cost required to increase agent complexity. At the
most basic level, replication of DNA molecules makes up the foundational
component of all biological life: the maintenance and development of
information, held up and carried forward against all odds of decay.

Understanding branching from a mathematical perspective has been a topic of
probability theory since the end of the nineteenth century. Darwin's cousin
Francis Galton sought to understand the persistence of family names over
generations, and for this, with Henry Watson, he employed what would come to be
known as the Galton--Watson process \cite{watson1875probability}. It is a
discrete-time stochastic process in which individuals either die or leave some
number of offspring at each generation, fate being dictated by defined
probabilities. In the binary case used throughout this thesis, each individual
present either replicates with probability $p$ or dies with the remaining
likelihood, so that the number $Z$ of offspring left by one individual has mean
\[
\mean{Z} = 2p + 0\,(1-p) = 2p .
\]

It wasn't until the 1920s that continuous time models of branching were
developed. In 1924, George Udny Yule published a paper in response to a book on
evolution. The book was called ``Age and area'' \cite{Willis1922}, and sought to understand the
relationship between the number of species in a living genus and the area over
which it spans. The paper \cite{yule1925ii} introduced what would become known
as the Yule process. This was the first continuous time model of branching.
Individuals only replicate, and do so with some probability during an interval
of time:
$$\pr{\xx_{t + \Delta t}  = \xt + 1 \;\; | \;\; \xt} = \lambda \xt \Delta t + o(\Delta t).$$
As the interval is taken in the limit, $\Delta t\to 0$, the probability that
more than one replication takes place becomes infinitesimal. The fact that any
of the $\xt$ present individuals can reproduce is encoded by the
proportionality of the birth rate $\lambda \xt$.

This model, applied equally well to bacterial replication as to the emergence
of species, is also known as the standard birth process. Introducing death
events, we get the standard birth--death process:

$$\pr{\xx_{t + \Delta t}  = \xt + 1 \;\; | \;\; \xt} = \lambda \xt \Delta t + o(\Delta t)$$
$$\pr{\xx_{t + \Delta t}  = \xt - 1 \;\; | \;\; \xt} = \mu \xt \Delta t + o(\Delta t),$$

\noindent
which now has a universally accessible absorbing state at $0$ representing
extinction.

\subsubsection*{Criticality}

One of the most important aspects of branching processes is their net
replacement, which asks whether an individual replaces itself on average. In
discrete time this is the expected number of offspring produced by a single
individual over its lifetime, $\mean{Z}=2p$, and the process is called
\textbf{supercritical}, \textbf{critical} or \textbf{subcritical} according as
$2p$ exceeds, equals or falls below $1$. In continuous time the same question
is asked of the mean population size, which obeys
$$\ddt{\mean{\xt}} = (\lambda - \mu) \mean{\xt},$$
so that the three regimes are set instead by the sign of the net replication
rate $\lambda-\mu$. The two conditions say the same thing in the two settings,
and the distinction between them is one of arithmetic rather than of substance.
The three regimes behave quite differently. Some supercritical realisations die
out and the rest rise towards infinity in an increasingly predictable
exponential fashion. Every subcritical realisation is destined to fixate at
$0$. So is every critical one --- but the expected time it takes to do so is
infinite, so that a single realisation may last a very long time without ever
fixating.

\subsubsection*{Catastrophes}

Much of the work presented in this thesis surrounds our analysis of a further
modification to this birth--death process. An additional ``catastrophe event''
is introduced, whereby a sudden transition can take place from any non-zero
count either to $0$ or to a separate ruptured state $H$, as per choice of
convention:

$$\pr{\xx_{t + \Delta t}  = \xt + 1 \;\; | \;\; \xt} = \lambda \xt \Delta t + o(\Delta t)$$
$$\pr{\xx_{t + \Delta t}  = \xt - 1 \;\; | \;\; \xt} = \mu \xt \Delta t + o(\Delta t)$$
$$\pr{\xx_{t + \Delta t}  = H \;\; | \;\; \xt} = \delta \xt \Delta t + o(\Delta t).$$

\noindent
Like birth and death, the catastrophe can be initiated by any of the individual
units, and so occurs at a rate proportional to the population size. In our case
this is to represent the possibility of any resident bacterium within a white
blood cell triggering a bursting event in which the walls of the macrophage
disintegrate and the pathogenic contents are simultaneously released.

The birth--death--catastrophe process is the spine of this thesis.
\Mextref{bdc:ch:core}{Chapter~4} and
\mextref{dist:ch:distribution-theory}{Chapter~5} solve the single cell;
\mextref{p:ch:one-cell-to-population}{Chapter~6} carries it up to a population
of cells; \mextref{tt:ch:two-type}{Chapter~7} solves a two-type version of it;
and \mextref{path:ch:pathogenesis}{Chapter~8} asks what the burst is worth
against a defence that can be used up.
